\documentclass[11pt]{article}

% 基础设置
\setlength{\parindent}{0pt} % 取消首行缩进
\usepackage{fontspec} % XeLaTeX核心字体包
\usepackage{xeCJK} % 中文支持
\CJKsetecglue{} % 取消中文与数字/字母间的间隙

% 字体设置 (请确保fonts文件夹路径正确，或使用系统字体)
\setmainfont{TeX Gyre Termes} % 英文衬线字体 (若无本地文件，直接调用系统名)
\setCJKmainfont{Noto Sans SC} % 中文无衬线字体

% 页面布局
\usepackage[
    a4paper,
    left=1.2cm,
    right=1.2cm,
    top=1.5cm,
    bottom=1cm,
    nohead
]{geometry}

% 颜色与链接
\usepackage{xcolor}
\definecolor{CVBlue}{RGB}{23,110,191}
\usepackage{hyperref}
\hypersetup{hidelinks} % 隐藏链接红色框
\usepackage{url}
\urlstyle{tt}

% 图标与图形
\usepackage{fontawesome} % 简历图标
\usepackage{tikz} % 用于定位 (如需头像)
\usetikzlibrary{calc}

% 标题与列表格式
\usepackage{titlesec}
\usepackage{enumitem}
\renewcommand{\baselinestretch}{1.3} % 行间距

% 自定义列表间距
\setlist{noitemsep, topsep=2pt} % 全局取消列表额外间距
\setlist[itemize]{leftmargin=*, label=\textbullet} % 自定义项目符号

% 自定义命令
% 标题格式：蓝色下划线
\titleformat{\section}
{\large\bfseries\raggedright}
{}{0em}
{}[{\color{CVBlue}\titlerule[0.8pt]}]
\titlespacing*{\section}{0cm}{*1.2}{*0.8}

% 项目描述与列表间的微小间距
\newcommand{\descsep}{\vspace{0.1em}}

\begin{document}
\pagenumbering{gobble} % 不显示页码

% --- 头像定位 (可选) ---
% 如果需要显示头像，请取消注释并确保图片存在
% \begin{tikzpicture}[remember picture, overlay]
%     \node[anchor=north east] at ($(current page.north east)+(-2cm,-1.2cm)$) {\includegraphics[height=3cm]{avatar.jpg}};
% \end{tikzpicture}%

% --- 个人信息 ---
\centerline{\LARGE\bfseries 王鑫}
\vspace{0.5em}
\centerline{\normalsize \faPhone\ 18512385345 \quad \faEnvelopeO\ \href{mailto:22515026@zju.edu.cn}{22515026@zju.edu.cn}}
\centerline{\normalsize \textcolor{CVBlue}{\textbf{目标岗位：Web Agent开发工程师}}}
\vspace{0.8em}

% --- 教育背景 ---
\section{\faGraduationCap\ 教育背景}
\textbf{浙江大学} \hfill 2025.9 -- 至今 \\
\textit{电子信息技术与仪器} \hfill 硕士 \\
\begin{itemize}
    \item \textbf{相关课程：} 《现代信号处理》、《人工神经网络与模式识别》、《模式识别与机器学习》
\end{itemize}

\textbf{浙江大学} \hfill 2021.9 -- 2025.6 \\
\textit{生物医学工程} \hfill 学士 \\
\begin{itemize}
    \item \textbf{相关课程：} 《C语言程序设计》、《高级程序设计》、《网络技术与应用》
\end{itemize}

% --- 项目经历 ---
\section{\faUsers\ 项目经历}

\textbf{基于字节跳动Eino框架的企业级 RAG Agent 系统} \hfill 2025.11 -- 2026.2
\begin{itemize}
    \item \textbf{项目描述：} 针对企业级Agent高并发与多步推理需求，基于字节跳动 \textbf{Eino} 框架与 \textbf{ADK}，整合 \textbf{Hertz} (Go Web框架) 与 \textbf{Milvus} (向量库) 开发生产级 RAG Agent。
    \item \textbf{核心开发：} 设计 \textbf{ReAct} 风格推理流程，将检索逻辑封装为 \textbf{Tool} 实现“思考-检索-生成”闭环；利用 Go 的 \textbf{goroutine} 优化编排层，解决异步任务阻塞问题。
    \item \textbf{项目成果：} 系统成功上线，支持 \textbf{数十路并发} 请求处理，沉淀出可复用的后端 Agent 模板，支撑后续浏览器自动化等场景。
\end{itemize}

\textbf{基于 Dify 的 RAG 智能问答系统后端设计} \hfill 2025.7 -- 2025.10
\begin{itemize}
    \item \textbf{架构设计：} 针对知识库查询效率痛点，基于 \textbf{Flask+Celery} 构建分离式架构，实现 \textbf{Agent+RAG} 协同工作流，支持 \textbf{多轮对话记忆} 与 \textbf{文档溯源}。
    \item \textbf{工程落地：} 使用 \textbf{Docker Compose} 一键编排 PostgreSQL、Redis、Weaviate 等中间件，解决环境依赖难题，实现从开发到生产的平滑迁移。
    \item \textbf{性能优化：} 优化动态检索策略，系统支持千份文档实时问答，平均响应时间控制在 \textbf{1.5秒} 以内。
\end{itemize}

% --- 技术栈 (横向排版，节省空间) ---
\section{\faCogs\ 技术栈}
\begin{tabular}{@{}ll@{}}
\textbf{语言:} & \textbf{Go} (高并发/协程) \quad \textbf{Python} (LLM/RAG) \quad \textbf{C++/Java} \\
\textbf{框架:} & \textbf{Eino/ADK} (Agent开发) \quad \textbf{Hertz/Flask} (Web服务) \quad \textbf{Playwright} (浏览器自动化) \\
\textbf{工具:} & \textbf{Docker Compose} (编排) \quad \textbf{Milvus/Weaviate} (向量库) \quad \textbf{Git/Vim} \\
\end{tabular}

% --- 外语与竞赛 ---
\section{\faInfo\ 外语与竞赛}
\begin{tabular}{@{}ll@{}}
\textbf{英语:} & CET-4 \textbf{616} \quad CET-6 \textbf{473} \\
\textbf{奖项:} & 全国大学生数学竞赛省一等奖 \\
\end{tabular}

\end{document}
